% page 634 du fac-simile (image p-633 du scan)
% bloc 154.9 x 246.3 mm ; marge gauche 8.0 mm
\pgc{-12.700mm}{0.000mm}{
\l{\marge[80.182mm]{145.661mm}{\ornamento{11.063mm}{ornements/letroj/litero-Z.png}}\cel{2}626}
\l{}
\l{}
\l{\sou{yugo}. I. Ligno-peco quan onu muntas sur la kapo di la bovi}
\l{\cel{3}e per qua onu jungas li duope. - II. (\sou{metaf.}) (a) Sub-}
\l{\cel{3}miseso quan impozas mastro. (b) Koakteso quan impozas}
\l{\cel{3}engajo, pasiono, e c. - III. (en\cel{1}\sou{la Roma antiqua}). Pi-}
\l{\cel{3}quo pozita horizontale sur du piqui vertikala, stekita}
\l{\cel{3}en sulo, sub qua onu igis pasar la enemiki vinkita por}
\l{\cel{3}humiligar e shamigar li. - DEFIS.}
\l{}
\l{\sou{yuko}. (\sou{bot.})\cel{2}Planto perena, ek la familio \textquotedbl{}lábiacei\textquotedbl{}, kun}
\l{\cel{3}folii longa, streta e harda, qui finas per pinto akuta ;}
\l{\cel{3}kun flori blanka, dispozita panikule sur shafto longa.}
\l{\cel{3}- L.\cel{1}\sou{yucca}.\cel{2}DEFIS.}
\l{}
\l{\sou{-yuno}.\cel{2}Sufixo qua indikas la filio tre yuna di animalo.}
\l{}
\l{\sou{yuna}.\cel{2}Qua ne evas multe. DEFI.}
\l{}
\l{\sou{yuro}. I. Lego-justifikeso pri postular ke onu retrodonez to}
\l{\cel{3}quon onu debas a ni. - II. Ensemblo ek la legi qui deter-}
\l{\cel{3}minas to quonsinglu, segun sua situeso, darfas legitime}
\l{\cel{3}agar (o ne-agar) o facar (o ne-facar) o postular de altru.}
\l{\cel{3}- DEFIRS.}
\l{}
\l{\sou{yusta}. Konforma a la yuro ; qua respektas la yuro di la}
\l{\cel{3}kun-homi. - EF.}
\l{}
\l{}
\l{}
\l{}
\l{}
\l{\cel{33}Z}
\l{}
\l{}
\l{}
\l{}
\l{\sou{zagayo}. Sorto di javelino, facata ed uzata da ula populi}
\l{\cel{3}sovaja. -}
\l{}
\l{\sou{zebro}. (\sou{zool.}) Animalo di Afrika, aspekto-vicina ye asno,}
\l{\cel{3}kun pilaro flava e strii bruna. L.\cel{1}\sou{equus zebra}. -DEFIRS}
\l{}
\l{\sou{zebuo}. (\sou{zool.)}\cel{2}Bovo kun un gibo, devenanta de India. -}
\l{\cel{3}L.\cel{1}\sou{bos taurus indicus}. - DEFIRS.}
\l{}
\l{\sou{zefiro.}\cel{2}I. (en\cel{1}\sou{la epoki antiqua}). Vento de la ocidento}
\l{\cel{3}equinoxala. - II. Irga vento lejera e dolca. - DEFIS.}
\l{}
\l{\sou{zeko}. (\sou{zool.})\cel{2}Insekto sen-ala qua su aplikas an la pelo}
\l{\cel{3}di la bovi, mutoni, hundi, e sugas lia sango. - L.}
\l{\cel{3}\sou{ixodes}. - DI.}
\l{}
\l{\sou{zelar}. (\sou{netrans.}) (\sou{por, pri}). Agar propra-move, sen tard-}
\l{\cel{3}eso, kun komplezo e fervoro, por servar (o +serviar)}
\l{\cel{3}ulu, ulo. - EFIS.}
}
